\section{Conclusion}\label{sec:conclusion}
We proposed C2D2, a two-pass training scheme that adds a synthetically degraded pass, distilled from a frozen clean-input teacher, to ordinary training. Applied to a fusion model on a frozen DINOv3 backbone with sensor-specific adapters, cross-modal attention, and a high-resolution path, it gives the highest test mIoU among published four-sensor fusion models on DELIVER while keeping the model robust when sensors are removed or corrupted, and it reaches 80.01 mIoU on the MUSES test server and 59.78 on MCubeS, above the best published four-sensor and overall results, respectively. Our analysis indicates that the heavy dependence of fusion models on depth in DELIVER comes mainly from training without degraded sensors \todo{wording after FU-04 and FU-05}. It also exposes what remains: the high-resolution path recovers thin objects but lets fine detail mislead the per-pixel classifier, as in posters read as traffic lights, which suggests routing fine detail to mask prediction rather than to per-pixel classification. Measuring the same dependence on real depth sensors is the next step that our synthetic-data findings call for.
